%============================ ANEXOS ========================================
\appendix
\chapter{Glosario de términos}
\label{anexo:glosario}
\begin{description}[leftmargin=2.2cm,style=nextline]
  \item[BLE] Bluetooth Low Energy; variante de bajo consumo del estándar Bluetooth.
  \item[BSSID] Identificador único (dirección MAC) de un punto de acceso Wi\nobreakdash-Fi.
  \item[Evil twin] Punto de acceso que suplanta el SSID de una red legítima.
  \item[Heatmap] Mapa de calor; representación de la ocupación por canal.
  \item[OUI] Organizationally Unique Identifier; prefijo de fabricante en la MAC.
  \item[Portal cautivo] Página que intercepta la navegación de un cliente recién conectado.
  \item[RSSI] Received Signal Strength Indicator; indicador de intensidad de señal.
  \item[SoftAP] Punto de acceso Wi\nobreakdash-Fi emitido por software desde el propio dispositivo.
  \item[SSID] Nombre público de una red Wi\nobreakdash-Fi.
  \item[WebSocket] Protocolo de comunicación bidireccional persistente sobre TCP.
  \item[WSS] WebSocket sobre TLS (cifrado).
\end{description}

\chapter{Referencias del código fuente}
\label{anexo:codigo}
El código fuente completo del sistema acompaña a este documento. La Tabla~\ref{tab:fuentes} enumera los archivos principales y su extensión.

\begin{table}[H]
\centering
\caption{Archivos de código fuente del proyecto}
\label{tab:fuentes}
\footnotesize
\begin{tabularx}{\linewidth}{@{}l r X@{}}
\toprule
\textbf{Archivo} & \textbf{Líneas} & \textbf{Descripción} \\
\midrule
\code{firmware/eilor\_esp32.ino} & 777 & Firmware del ESP32 (4 modos, portal cautivo, buffer) \\
\code{server.js} & 665 & Servidor Express + WebSocket + API REST \\
\code{lib/analysis.js} & 71 & Puntaje de canal y detección de gemelo malicioso \\
\code{lib/oui.js} & 62 & Fabricante por OUI y MAC aleatorias \\
\code{lib/store.js} & 81 & Persistencia JSON y series temporales \\
\code{lib/auth.js} & 60 & Sesiones por token y autorización \\
\code{public/app.js} & 913 & Lógica del dashboard en tiempo real \\
\code{public/index.html} & 419 & Estructura del dashboard \\
\code{public/styles.css} & 1272 & Estilos del tablero \\
\bottomrule
\end{tabularx}
\end{table}

\chapter{Repositorio y despliegue de referencia}
\label{anexo:repositorio}

El código fuente completo del proyecto (firmware, servidor, módulos y frontend) se
aloja en el siguiente repositorio, que acompaña a este documento como anexo digital:

\begin{itemize}[leftmargin=1.4em]
  \item \textbf{Repositorio de código:} \repourl
  \item \textbf{Plataforma en producción:} \plataformaurl
  \item \textbf{Túnel y hostname público:} \code{esp.orellanadigitalco.trade} (puerto 443, WSS/HTTPS), publicado con Cloudflare Tunnel (Capítulo~\ref{cap:cloudflare}).
\end{itemize}

\noindent La plataforma en producción integra el flujo de extremo a extremo
descrito en el Capítulo~\ref{cap:plataforma}: el dispositivo ESP32 se conecta por
WSS al servidor a través del túnel, y los clientes web consumen el \emph{dashboard}
sobre HTTPS. Para reproducir el despliegue basta clonar el repositorio, instalar las
dependencias (\code{npm install}), definir \code{EILOR\_PASSWORD} y ejecutar el
servidor bajo PM2, siguiendo el manual técnico del Capítulo~\ref{cap:manuales}.
